% !TeX root = ../../Book1.tex
% 纲要标题：### 10.4 中国剩余定理
% 纲要位置：计划纲要.md，第 406 行。

\section{中国剩余定理}
\label{b1:sec:1004}

整数的两个互素模数可以独立指定余数。对一般交换环，替代“模数互素”的条件是两个理想的和等于整个环。中国剩余定理不仅给出同余方程的解，还把一个商环拆成若干分量；逆向看，这些分量由特殊的幂等元识别。本节除最后的辨析外均在交换含幺环中讨论。

% 纲要内容（仅作写作计划，尚非教材正文）：
%
% * 两两互素理想与有限直积分解
% * 整数同余方程
% * 幂等元与环的分解
%
% ---
%

\subsection{两两互素理想与有限直积分解}
\label{b1:subsec:1004:01}
\begin{definition}\label{b1:def:comaximal-ideals}
设 \(R\) 为交换环，\(I,J\) 为 \(R\) 的理想。若 \(I+J=R\)，则称 \(I,J\) 为\term[husu-lixiang]{互素理想}[comaximal ideals]。这等价于存在 \(a\in I,b\in J\) 使 \(a+b=1\)。
\end{definition}

\begin{lemma}\label{b1:lem:comaximal-product}
设 \(I,J\) 为交换环 \(R\) 的互素理想，则 \(I\cap J=IJ\)。更一般地，若 \(r\) 为正整数，且 \(R\) 的理想 \(I_1,\ldots,I_r\) 两两互素，则
\(\bigcap_{i=1}^r I_i=\prod_{i=1}^r I_i\)。
\end{lemma}
\begin{proof}
\cref{b1:prop:ideal-containments}已给出 \(IJ\subseteq I\cap J\)。若 \(x\in I\cap J\)，取 \(a+b=1\)、\(a\in I,b\in J\)，则 \(x=xa+xb\in IJ\)，这里交换性把 \(xa\) 写成 \(ax\)。

对多个理想，取 \(a_j\in I_j,b_j\in I_r\) 使 \(a_j+b_j=1\)，其中 \(j<r\)。乘积 \(a_1\cdots a_{r-1}\) 属于 \(I_1\cdots I_{r-1}\)，且模 \(I_r\) 同余于 \(1\)。所以 \(I_r+I_1\cdots I_{r-1}=R\)。从 \(r=2\) 起归纳，将两理想结论用于这个乘积理想与 \(I_r\)，即可得到有限情形。
\end{proof}
\term[zhongguo-shengyu-dingli]{中国剩余定理}[Chinese remainder theorem]把这一等式与一个满同态结合起来。

\begin{theorem}[中国剩余定理]\label{b1:thm:crt}
设 \(r\geq1\) 为整数，交换环 \(R\) 的理想 \(I_1,\ldots,I_r\) 两两互素。则映射
\[
\Phi:R\longrightarrow\prod_{i=1}^rR/I_i,\qquad
a\longmapsto(a+I_i)_{i=1}^r
\]
是满环同态，核为 \(\bigcap_iI_i=\prod_iI_i\)，因而诱导环同构
\[
R/\Bigl(\prod_{i=1}^rI_i\Bigr)\cong\prod_{i=1}^rR/I_i.
\]
反过来，若 \(\Phi\) 满射，则这些理想必两两互素。
\end{theorem}
\begin{proof}
逐坐标考察可知 \(\Phi\) 保持运算与单位，其核恰为全部 \(I_i\) 的交。对每个 \(i\ne j\)，取 \(u_{ij}\in I_i,v_{ij}\in I_j\) 使 \(u_{ij}+v_{ij}=1\)，并令
\[
e_i=\prod_{j\ne i}v_{ij}.
\]
空乘积为 \(1\)。于是 \(e_i\equiv1\pmod{I_i}\)，且 \(e_i\in I_j\) 对 \(j\ne i\) 成立。给定任意余数组 \((a_i+I_i)_i\)，元素 \(a=\sum_i a_ie_i\) 的第 \(i\) 个余数恰为 \(a_i+I_i\)，所以 \(\Phi\) 满射。由\cref{b1:lem:comaximal-product}确定核，再对 \(\Phi\) 应用\cref{b1:thm:ring-first-iso}，得到所述同构。

反过来，若 \(\Phi\) 满射，对 \(i\ne j\) 选择一个原像 \(a\)，使 \(a\equiv1\pmod{I_i}\)、\(a\equiv0\pmod{I_j}\)。则 \(1=(1-a)+a\in I_i+I_j\)，故两理想互素。
\end{proof}
有限性使构造中的求和和乘积都在环内有意义。不能直接把定理中的有限直积改成无限直积；例如 \(\mathbb Z\) 只有可数多个元素，而对全部素数 \(p\) 的直积 \(\prod_p\mathbb F_p\) 不可数，自然映射不可能满射。

\subsection{整数同余方程}
\label{b1:subsec:1004:02}
正整数 \(m,n\) 满足 \(m\mathbb Z+n\mathbb Z=\mathbb Z\)，当且仅当 \(\gcd(m,n)=1\)。若 \(um+vn=1\)，方程
\[
x\equiv a\pmod m,\qquad x\equiv b\pmod n
\]
的解可取 \(x=avn+bum\)。模 \(m\) 第一项同余于 \(a\)、第二项为零；模 \(n\) 恰好相反。\cref{b1:thm:crt}说明全部解组成模 \(mn\) 的唯一剩余类。

\ProseParagraphBreak
例如求 \(x\equiv2\pmod3\)、\(x\equiv3\pmod4\)。由于 \(4-3=1\)，可以取 \(x=2\cdot4+3\cdot(-3)=-1\)，所以答案是 \(x\equiv11\pmod{12}\)。直接核验 \(11\) 的两个余数，可同时检查系数的放置次序。

\begin{proposition}\label{b1:prop:integer-crt-compatible}
设 \(m,n\) 为正整数，\(a,b\in\mathbb Z\)，并令 \(d=\gcd(m,n)\)。同余方程组
\[
x\equiv a\pmod m,\qquad x\equiv b\pmod n
\]
有整数解，当且仅当 \(a\equiv b\pmod d\)。有解时，全部解为模 \(\operatorname{lcm}(m,n)\) 的一个剩余类。
\end{proposition}
\begin{proof}
若有解，\(d\) 同时整除 \(x-a,x-b\)，所以整除 \(a-b\)。反过来，写 \(x=a+mt\)，所求条件成为 \(mt\equiv b-a\pmod n\)。若 \(d\mid b-a\)，约去 \(d\) 后，\(m/d\) 与 \(n/d\) 互素，Bézout 等式使 \(m/d\) 在模 \(n/d\) 下可逆，因而存在 \(t\)。两个解之差必须同时被 \(m,n\) 整除，恰好被它们的最小公倍数整除，且加上这个公倍数保持两个同余式，故得到全部解。
\end{proof}
若 \(n=p_1^{a_1}\cdots p_r^{a_r}\) 为整数的素数幂分解，则
\[
\mathbb Z/n\mathbb Z\cong\prod_{i=1}^r\mathbb Z/p_i^{a_i}\mathbb Z.
\]
这使大模数的乘法、可逆性与幂等元计算分别归结到各个素数幂分量。这里用的是整数的唯一分解；下一章才将元素分解的问题推广到一般整环。

\subsection{幂等元与环的分解}
\label{b1:subsec:1004:03}
\begin{definition}\label{b1:def:idempotent}
设 \(R\) 为环，\(e\in R\)。若 \(e^2=e\)，则称 \(e\) 为\term[midengyuan]{幂等元}[idempotent element]。若一族幂等元 \((e_i)_{i\in I}\) 对任意不同指标 \(i,j\in I\) 都满足 \(e_ie_j=e_je_i=0\)，则称这族幂等元两两正交。
\end{definition}

\begin{proposition}\label{b1:prop:idempotent-product}
设 \(r\geq1\) 为整数。若交换环 \(R\) 有两两正交的幂等元 \(e_1,\ldots,e_r\)，且 \(\sum_i e_i=1\)，则
\[
R\longrightarrow\prod_{i=1}^r e_iR,\qquad
a\longmapsto(e_ia)_i
\]
为环同构。每个 \(e_iR\) 的单位是 \(e_i\)，一般不等于 \(1_R\)。反之，若 \(R_1,\ldots,R_r\) 为交换环，并令
\[
R=\prod_{i=1}^rR_i,
\qquad
e_i=(0,\ldots,0,1_{R_i},0,\ldots,0),
\]
其中 \(1_{R_i}\) 位于第 \(i\) 个坐标，则 \(e_1,\ldots,e_r\) 两两正交、各自幂等，且 \(\sum_i e_i=1_R\)。
\end{proposition}
\begin{proof}
\(e_iR\) 对加法及乘法封闭，\(e_i(e_ia)=e_ia\)，故以 \(e_i\) 为单位。映射保持乘法，因为 \((e_ia)(e_ib)=e_iab\)，并把 \(1\) 送到各分量的单位。若每个 \(e_ia=0\)，则 \(a=\sum_i e_ia=0\)，所以单射。给定 \(x_i\in e_iR\)，令 \(a=\sum_i x_i\)；由正交性，\(e_ia=x_i\)，故满射。逆映射就是各分量求和。有限直积中，第 \(i\) 坐标为单位而其余为零的元素显然满足幂等、正交及和为单位的条件。
\end{proof}
当 \(e_i\ne1_R\) 时，包含映射 \(e_iR\hookrightarrow R\) 不保持单位，因而不是本书约定下的环同态，也不是本书约定下的子环嵌入；\(e_iR\) 在上述分解中以 \(e_i\) 为单位独立成环。

\ProseParagraphBreak
这些幂等元也直接给出中国剩余定理中的理想。对每个 \(i\)，映射
\[
p_i:R\longrightarrow e_iR,\qquad a\longmapsto e_i a
\]
是满环同态，并且 \(p_i(1_R)=e_i=1_{e_iR}\)，所以保持单位。若 \(e_i a=0\)，则 \(a=(1-e_i)a\)；反过来，\(e_i(1-e_i)=0\)。因此
\[
J_i:=\ker p_i=(1-e_i)R,
\qquad e_iR\cong R/J_i.
\]
对于 \(i\ne j\)，由 \(1=(1-e_i)+e_i\) 和 \(e_i=(1-e_j)e_i\in J_j\) 可知 \(J_i+J_j=R\)。若 \(a\in\bigcap_iJ_i\)，则每个 \(e_i a=0\)，于是 \(a=\sum_i e_i a=0\)。故 \(\bigcap_iJ_i=0\)，再由\cref{b1:thm:crt}得到
\[
R\cong\prod_{i=1}^rR/J_i\cong\prod_{i=1}^r e_iR.
\]
这个同构在第 \(i\) 个坐标上的映射就是 \(p_i\)。反过来，从两两互素的理想出发时，\cref{b1:thm:crt}证明中构造的元素 \(e_i\) 未必在原环中幂等，但它们在商环 \(R/\bigcap I_i\) 中的像就是这些坐标幂等元。例如 \(\mathbb Z/12\mathbb Z\cong\mathbb Z/3\mathbb Z\times\mathbb Z/4\mathbb Z\) 中，\([4]\) 与 \([9]\) 分别对应 \((1,0)\)、\((0,1)\)：它们的平方各等于自身，乘积为零，和为 \([1]\)。

\ProseParagraphBreak
整环中 \(e^2=e\) 等价于 \(e(e-1)=0\)，所以只有 \(0,1\) 两个幂等元，不能分成两个非零环的直积。反方向不成立：\(\mathbb Z/4\mathbb Z\) 含零因子，却只有 \([0],[1]\) 两个幂等元。一般非交换环要作上述直积分解，还应要求幂等元与所有元素交换，即为中心幂等元。这样证明中的 \((e_ia)(e_ib)=e_iab\) 才成立；仅有一个矩阵幂等元 \(E_{11}\) 并不提供矩阵环的直积分解。
